\section{并发编程}

并发程序是由在时间上重叠的一组逻辑流组成的。现代操作系统提供了三种基本的构造并发程序的方法：
\begin{itemize}
    \item \textbf{进程}：每个逻辑控制流都是一个进程，由内核调度和维护。进程有独立的虚拟地址空间，需要显式的进程间通信（IPC）机制来共享状态。
    \item \textbf{I/O多路复用}：应用程序在一个进程上下文中显式地调度逻辑流。逻辑流被模型化为状态机，所有流共享同一个地址空间。
    \item \textbf{线程}：运行在单一进程上下文中的逻辑流，由内核调度。线程像进程流一样由内核调度，像I/O多路复用流一样共享同一个虚拟地址空间。
\end{itemize}

应用级并发在以下场景中很有用：
\begin{itemize}
    \item 访问慢速I/O设备（磁盘等）
    \item 与人交互（窗口系统）
    \item 通过推迟工作以降低延迟
    \item 服务多个网络客户端
    \item 在多核机器上进行并行计算
\end{itemize}

\subsection{基于进程的并发编程}

\subsubsection{原理}
构造并发程序最简单的方法是使用进程。父进程接受客户端连接请求，然后创建子进程为每个新客户端提供服务。因为父、子进程中的已连接描述符都指向同一个文件表表项，所以父进程关闭它的已连接描述符副本是至关重要的，否则会导致内存泄漏。

\subsubsection{基于进程的并发服务器}
\begin{lstlisting}[style=CStyle]
#include "csapp.h"

void echo(int connfd);
void sigchld_handler(int sig) {
    while (waitpid(-1, 0, WNOHANG) > 0);
    return;
}

int main(int argc, char **argv) {
    int listenfd, connfd;
    socklen_t clientlen;
    struct sockaddr_storage clientaddr;
    
    if (argc != 2) {
        fprintf(stderr, "usage: %s <port>\n", argv[0]);
        exit(0);
    }
    
    Signal(SIGCHLD, sigchld_handler);
    listenfd = Open_listenfd(argv[1]);
    
    while (1) {
        clientlen = sizeof(struct sockaddr_storage);
        connfd = Accept(listenfd, (SA*)&clientaddr, &clientlen);
        
        if (Fork() == 0) { // 子进程
            Close(listenfd); // 子进程关闭监听描述符
            echo(connfd);    // 服务客户端
            Close(connfd);   // 关闭连接
            exit(0);
        }
        Close(connfd); // 父进程关闭已连接描述符
    }
}
\end{lstlisting}

\subsubsection{进程的优劣}
\textbf{优点}：
\begin{itemize}
    \item 进程有独立的地址空间，一个进程不可能不小心覆盖另一个进程的虚拟内存
    \item 消除了许多令人迷惑的错误
\end{itemize}

\textbf{缺点}：
\begin{itemize}
    \item 共享状态信息困难，需要显式的IPC机制
    \item 进程控制开销高，性能较差
\end{itemize}

\subsection{基于I/O多路复用的并发编程}

\subsubsection{I/O多路复用原理}
I/O多路复用技术使用\texttt{select}函数，要求内核挂起进程，只有在一个或多个I/O事件发生后，才将控制返回给应用程序。基本思路是：
\begin{itemize}
    \item \texttt{select}等待一组描述符准备好读
    \item 当描述符准备好时，内核返回控制给应用程序
    \item 应用程序使用宏操作描述符集合：\texttt{FD\_ZERO}, \texttt{FD\_SET}, \texttt{FD\_CLR}, \texttt{FD\_ISSET}
\end{itemize}

\subsubsection{事件驱动编程}
在事件驱动程序中，逻辑流被模型化为状态机。一个状态机由一组状态、输入事件和转移组成，其中转移是将状态和输入事件映射到状态。

\subsubsection{基于I/O多路复用的并发事件驱动服务器}
\begin{lstlisting}[style=CStyle]
#include "csapp.h"

typedef struct { // 客户端池
    int maxfd;           // 读集合中最大描述符
    fd_set read_set;     // 所有活动描述符集合
    fd_set ready_set;    // 准备好读的描述符子集
    int nready;          // select返回的准备好数
    int maxi;            // clientfd数组最大索引
    int clientfd[FD_SETSIZE]; // 活动描述符集合
    rio_t clientrio[FD_SETSIZE]; // 活动读缓冲区
} pool;

int byte_cnt = 0; // 服务器接收字节总数

void init_pool(int listenfd, pool *p) {
    // 初始化代码
}

void add_client(int connfd, pool *p) {
    // 添加客户端代码
}

void check_clients(pool *p) {
    // 服务客户端代码
}

int main(int argc, char **argv) {
    int listenfd, connfd;
    socklen_t clientlen;
    struct sockaddr_storage clientaddr;
    static pool pool;
    
    if (argc != 2) {
        fprintf(stderr, "usage: %s <port>\n", argv[0]);
        exit(0);
    }
    
    listenfd = Open_listenfd(argv[1]);
    init_pool(listenfd, &pool);
    
    while (1) {
        // 等待描述符准备就绪
        pool.ready_set = pool.read_set;
        pool.nready = Select(pool.maxfd+1, &pool.ready_set, 
                             NULL, NULL, NULL);
        
        // 处理新连接
        if (FD_ISSET(listenfd, &pool.ready_set)) {
            clientlen = sizeof(struct sockaddr_storage);
            connfd = Accept(listenfd, (SA*)&clientaddr, &clientlen);
            add_client(connfd, &pool);
        }
            
        // 服务准备好的客户端
        check_clients(&pool);
    }
}
\end{lstlisting}

\subsubsection{I/O多路复用技术的优劣}
\textbf{优点}：
\begin{itemize}
    \item 程序员对程序行为有更多控制
    \item 所有流共享进程地址空间，便于数据共享
    \item 可使用熟悉调试工具（如GDB）
    \item 通常比基于进程的设计更高效
\end{itemize}

\textbf{缺点}：
\begin{itemize}
    \item 编码复杂，随着并发粒度减小复杂度上升
    \item 不能充分利用多核处理器
    \item 对恶意客户端攻击（如只发送部分文本行）较为脆弱
\end{itemize}

\subsection{基于线程的并发编程}

\subsubsection{线程基础}
线程是运行在进程上下文中的逻辑流，由内核自动调度。每个线程有自己的线程上下文：
\begin{itemize}
    \item 唯一整数线程ID
    \item 栈、栈指针、程序计数器
    \item 通用目的寄存器、条件码
\end{itemize}

所有运行在一个进程里的线程共享该进程的整个虚拟地址空间，包括代码、数据、堆、共享库和打开的文件。

\subsubsection{线程执行模型}
\begin{itemize}
    \item 每个进程开始时都是单一主线程
    \item 主线程创建对等线程，之后并发运行
    \item 线程上下文切换比进程上下文切换快得多
    \item 线程不是按严格父子层次组织，而是组成对等线程池
    \item 任意线程可以杀死或等待任意对等线程终止
\end{itemize}

\subsubsection{Pthreads}
POSIX线程（Pthreads）是处理线程的标准接口。关键函数包括：
\begin{itemize}
    \item \texttt{pthread\_create}：创建新线程
    \item \texttt{pthread\_exit}：显式终止线程
    \item \texttt{pthread\_join}：等待线程终止
    \item \texttt{pthread\_detach}：分离线程
    \item \texttt{pthread\_once}：初始化线程
    \item \texttt{pthread\_cancel}：取消线程
\end{itemize}

\subsubsection{基于线程的并发服务器}
\begin{lstlisting}[style=CStyle]
#include "csapp.h"

void echo(int connfd);
void *thread(void *vargp);

int main(int argc, char **argv) {
    int listenfd, *connfdp;
    socklen_t clientlen;
    struct sockaddr_storage clientaddr;
    pthread_t tid;
    
    if (argc != 2) {
        fprintf(stderr, "usage: %s <port>\n", argv[0]);
        exit(0);
    }
    
    listenfd = Open_listenfd(argv[1]);
    
    while (1) {
        clientlen = sizeof(struct sockaddr_storage);
        connfdp = Malloc(sizeof(int));
        *connfdp = Accept(listenfd, (SA*)&clientaddr, &clientlen);
        
        // 创建对等线程处理请求
        Pthread_create(&tid, NULL, thread, connfdp);
    }
}

void *thread(void *vargp) {
    int connfd = *((int *)vargp);
    Pthread_detach(pthread_self()); // 分离线程
    Free(vargp);                    // 释放内存
    
    echo(connfd);                   // 服务客户端
    Close(connfd);                  // 关闭连接
    return NULL;
}
\end{lstlisting}

\subsection{多线程程序中的共享变量}

\subsubsection{线程内存模型}
线程的内存模型包含以下关键点：
\begin{itemize}
    \item 寄存器：从不共享
    \item 虚拟内存：总是共享（整个用户虚拟地址空间）
    \item 线程栈：通常独立访问，但不设防（如果一个线程获取了指向其他线程栈的指针，可以读写该栈）
\end{itemize}

\subsubsection{变量映射到内存}
\begin{itemize}
    \item \textbf{全局变量}：在运行时，虚拟内存的读/写区域只包含每个全局变量的一个实例，任何线程都可以引用。
    \item \textbf{本地自动变量}：每个线程的栈都包含它自己的所有本地自动变量的实例。
    \item \textbf{本地静态变量}：虚拟内存的读/写区域只包含每个本地静态变量的一个实例，所有线程共享。
\end{itemize}

\subsubsection{共享变量}
一个变量是共享的，当且仅当它的一个实例被一个以上的线程引用。共享变量引入了同步错误的可能性。

\subsection{用信号量同步线程}

\subsubsection{同步问题}
当多个线程修改共享变量时，可能导致错误结果。例如，两个线程对共享计数器加1，可能得到错误结果：
\begin{itemize}
    \item 预期：$cnt = 2 \times niters$
    \item 实际：每次运行得到不同错误结果
\end{itemize}

进度图（Progress Graph）将$n$个并发线程的执行模型化为$n$维笛卡尔空间中的轨迹。关键概念：
\begin{itemize}
    \item \textbf{临界区}：操作共享变量的指令区域
    \item \textbf{不安全区}：临界区交集形成的区域
    \item \textbf{安全轨迹线}：绕开不安全区的轨迹
    \item \textbf{不安全轨迹线}：接触到不安全区的轨迹
\end{itemize}

\begin{figure}[H]
    \centering
    \includegraphics[width=0.8\textwidth]{CONC/progress-graph.png}
\end{figure}

\subsubsection{信号量}
信号量是具有非负整数值的全局变量，只能由两种特殊操作处理：
\begin{itemize}
    \item \texttt{P(s)}（Wait）：如果$s$非零，将$s$减1并返回；如果$s$为零，挂起线程直到$s$变为非零。
    \item \texttt{V(s)}（Signal）：将$s$加1。如果有线程阻塞在P操作等待$s$变为非零，重启其中一个。
\end{itemize}

POSIX信号量函数：
\begin{lstlisting}[language=C]
#include <semaphore.h>
int sem_init(sem_t *sem, 0, unsigned int value);
int sem_wait(sem_t *s);  // P(s)
int sem_post(sem_t *s);  // V(s)
\end{lstlisting}

\subsubsection{使用信号量实现互斥}
二元信号量（互斥锁）：值总是0或1。通过P和V操作包围临界区：
\begin{lstlisting}[language=C]
volatile long cnt = 0;
sem_t mutex;

// 初始化
Sem_init(&mutex, 0, 1);

// 线程函数
for (i = 0; i < niters; i++) {
    P(&mutex);
    cnt++;
    V(&mutex);
}
\end{lstlisting}

\begin{figure}[H]
    \centering
    \includegraphics[width=0.8\textwidth]{CONC/PV-graph.png}
\end{figure}

\subsubsection{使用信号量调度共享资源}

\textbf{生产者-消费者问题}：有限缓冲区，生产者插入项目，消费者取出项目。需要：
\begin{itemize}
    \item 互斥访问（mutex信号量）
    \item 调度访问（slots和items信号量）
\end{itemize}

SBUF: 同步有限缓冲区
\begin{lstlisting}[style=CStyle]
typedef struct {
    int *buf;        // 缓冲区数组
    int n;           // 最大槽位数
    int front;       // buf[(front+1)%n]是第一个项目
    int rear;        // buf[rear%n]是最后一个项目
    sem_t mutex;     // 保护对buf的访问
    sem_t slots;     // 可用槽位数
    sem_t items;     // 可用项目数
} sbuf_t;

void sbuf_insert(sbuf_t *sp, int item) {
    P(&sp->slots);    // 等待可用槽位
    P(&sp->mutex);    // 锁定缓冲区
    sp->buf[(++sp->rear)%(sp->n)] = item;
    V(&sp->mutex);    // 解锁缓冲区
    V(&sp->items);    // 通知有新项目
}

int sbuf_remove(sbuf_t *sp) {
    int item;
    P(&sp->items);    // 等待可用项目
    P(&sp->mutex);    // 锁定缓冲区
    item = sp->buf[(++sp->front)%(sp->n)];
    V(&sp->mutex);    // 解锁缓冲区
    V(&sp->slots);    // 通知有新槽位
    return item;
}
\end{lstlisting}

\textbf{读者-写者问题}：
\begin{itemize}
    \item 读者：只读共享对象
    \item 写者：修改共享对象
    \item 两类问题：
    \begin{itemize}
        \item 读者优先：不要让读者等待，除非已给写者权限
        \item 写者优先：一旦写者准备好，尽可能快完成
    \end{itemize}
\end{itemize}

\paragraph{读者优先解决方案}

关键思想：
\begin{itemize}
    \item 信号量$w$控制对共享对象的访问
    \item 信号量$mutex$保护共享变量$readcnt$
    \item $readcnt$统计当前在临界区中的读者数量
    \item 第一个读者获取$w$，最后一个读者释放$w$
    \item 写者对$w$加锁/解锁，确保独占访问
\end{itemize}

\begin{lstlisting}[style=CStyle]
#include <pthread.h>
#include <semaphore.h>

// 全局共享变量
int readcnt = 0;       // 当前读者数量
sem_t mutex, w;        // 信号量，初始值都为1

void init_rw() {
    sem_init(&mutex, 0, 1);  // 保护readcnt
    sem_init(&w, 0, 1);      // 保护共享数据
}

void reader(void) {
    while (1) {
        P(&mutex);           // 进入临界区
        readcnt++;
        if (readcnt == 1)    // 第一个读者
            P(&w);           // 阻止写者
        V(&mutex);           // 离开临界区
        
        // 临界区：读操作
        printf("Reader %ld is reading\n", pthread_self());
        
        P(&mutex);           // 进入临界区
        readcnt--;
        if (readcnt == 0)    // 最后一个读者
            V(&w);           // 允许写者
        V(&mutex);           // 离开临界区
    }
}

void writer(void) {
    while (1) {
        P(&w);               // 进入临界区
        
        // 临界区：写操作
        printf("Writer %ld is writing\n", pthread_self());
        
        V(&w);               // 离开临界区
    }
}
\end{lstlisting}

\paragraph{写者优先解决方案}

写者优先的实现更复杂，需要额外的信号量来确保写者不会饥饿。

\begin{lstlisting}[style=CStyle]
#include <pthread.h>
#include <semaphore.h>

// 全局共享变量
int readcnt = 0;       // 当前读者数量
int writecnt = 0;      // 当前写者数量
sem_t rmutex, wmutex;  // 保护readcnt和writecnt
sem_t readTry;         // 读者尝试锁
sem_t resource;        // 保护共享资源

void init_writer_priority() {
    sem_init(&rmutex, 0, 1);   // 保护readcnt
    sem_init(&wmutex, 0, 1);   // 保护writecnt
    sem_init(&readTry, 0, 1);  // 读者尝试进入
    sem_init(&resource, 0, 1); // 保护共享资源
}

void reader_writer_priority(void) {
    while (1) {
        P(&readTry);           // 读者尝试进入
        P(&rmutex);            // 进入读者临界区
        readcnt++;
        if (readcnt == 1)
            P(&resource);      // 阻止写者
        V(&rmutex);            // 离开读者临界区
        V(&readTry);           // 允许后续读者进入
        
        // 临界区：读操作
        printf("Reader %ld is reading (writer priority)\n", pthread_self());
        
        P(&rmutex);            // 进入读者临界区
        readcnt--;
        if (readcnt == 0)
            V(&resource);      // 允许写者
        V(&rmutex);            // 离开读者临界区
    }
}

void writer_writer_priority(void) {
    while (1) {
        P(&wmutex);            // 进入写者临界区
        writecnt++;
        if (writecnt == 1)
            P(&readTry);       // 阻止后续读者
        V(&wmutex);            // 离开写者临界区
        
        P(&resource);          // 进入资源临界区
        
        // 临界区：写操作
        printf("Writer %ld is writing (writer priority)\n", pthread_self());
        
        V(&resource);          // 离开资源临界区
        
        P(&wmutex);            // 进入写者临界区
        writecnt--;
        if (writecnt == 0)
            V(&readTry);       // 允许读者
        V(&wmutex);            // 离开写者临界区
    }
}
\end{lstlisting}

\subsubsection{基于预线程化的并发服务器}
预线程化：使用固定数量的线程池处理请求，避免为每个请求创建线程的开销。
\begin{lstlisting}[style=CStyle]
#include "csapp.h"
#include "sbuf.h"

#define NTHREADS 4    // 工作者线程数
#define SBUFSIZE 16   // 缓冲区大小

sbuf_t sbuf;          // 共享连接描述符缓冲区

int main(int argc, char **argv) {
    // ...初始化
    
    sbuf_init(&sbuf, SBUFSIZE);
    
    // 创建工作者线程
    for (i = 0; i < NTHREADS; i++)
        Pthread_create(&tid, NULL, thread, NULL);
        
    while (1) {
        connfd = Accept(listenfd, (SA*)&clientaddr, &clientlen);
        sbuf_insert(&sbuf, connfd); // 插入缓冲区
    }
}

void *thread(void *vargp) {
    Pthread_detach(pthread_self());
    
    while (1) {
        int connfd = sbuf_remove(&sbuf); // 从缓冲区取出
        echo_cnt(connfd);                // 服务客户端
        Close(connfd);
    }
}
\end{lstlisting}

\subsection{其他并发问题}

\subsubsection{线程安全}
线程安全函数：被多个并发线程反复调用时，一直产生正确结果。

\textbf{线程不安全函数分类}：
\begin{enumerate}
    \item 不保护共享变量的函数：使用同步操作保护
    \item 保持跨越多个调用状态的函数：重写，不使用静态数据
    \item 返回指向静态变量指针的函数：重写或加锁-复制
    \item 调用线程不安全函数的函数：取决于被调用函数类型
\end{enumerate}

\subsubsection{可重入性}
可重入函数在被多个线程调用时，不引用任何共享数据。可重入函数是线程安全函数的真子集。

\textbf{显式可重入}：所有参数传值，数据引用为本地自动栈变量

\textbf{隐式可重入}：允许指针参数，但调用者需传递非共享数据指针

\subsubsection{常见线程不安全库函数}
\begin{table}[H]
\centering
\begin{tabular}{lll}
\toprule
线程不安全函数 & 线程不安全类 & 线程安全版本 \\
\midrule
rand & 2 & rand\_r \\
strtok & 2 & strtok\_r \\
asctime, ctime & 3 & asctime\_r, ctime\_r \\
gethostbyaddr, gethostbyname & 3 & gethostbyname\_r \\
inet\_ntoa & 3 & (无直接替代) \\
localtime & 3 & localtime\_r \\
\bottomrule
\end{tabular}
\end{table}

\subsubsection{竞争}
竞争：程序正确性依赖于线程特定执行顺序。

\textbf{竞争示例}：
\begin{lstlisting}[language=C]
// 问题代码：竞争条件
for (i = 0; i < N; i++)
    Pthread_create(&tid[i], NULL, thread, &i);

// 解决方案：为每个线程分配独立内存
for (i = 0; i < N; i++) {
    ptr = Malloc(sizeof(int));
    *ptr = i;
    Pthread_create(&tid[i], NULL, thread, ptr);
}
\end{lstlisting}

\subsubsection{死锁}
死锁：一组线程被阻塞，等待永远不为真的条件。

\begin{figure}[H]
    \centering
    \includegraphics[width=0.8\textwidth]{CONC/deadlock-graph.png}
\end{figure}

\textbf{互斥锁加锁顺序规则}：给定所有互斥操作的一个全序，如果每个线程都是以一种顺序获得互斥锁并以相反的顺序释放，那么这个程序就是无死锁的。

\newpage